\chapter{Machine learning for fault detection}
\label{ch:ml}

This chapter covers the data-driven diagnosis pipeline in \code{pyplant/ml}:
how runs become feature windows, the theory behind each classical detector,
how thresholds and the health index are set, and how a detection is traced
back to an instrument. The neural detector is in Chapter~\ref{ch:dl} and the
evaluation method in Chapter~\ref{ch:eval}.

\section{The pipeline}

\begin{figure}[htbp]
\centering
\scalebox{0.8}{%
\begin{tikzpicture}[node distance=6mm and 9mm]
  \node[box] (grid) {scenario grid\\\code{default_dataset()}};
  \node[box,right=of grid] (sim) {fast runner\\$\approx$1900$\times$ real time};
  \node[box,right=of sim] (pq) {\code{dataset.parquet}\\labelled samples};
  \node[box,below=of pq] (split) {split by scenario\\train / calibrate / test};
  \node[box,left=of split] (win) {windows + features\\$12\times4=48$ columns};
  \node[hot,left=of win] (fit) {fit on normal\\training runs};
  \node[box,below=of fit] (cal) {threshold = quantile\\on calibration runs};
  \node[box,right=of cal] (ev) {evaluate on test runs\\FAR, delay, isolation};
  \node[good,right=of ev] (sel) {select best\\\code{detector.joblib}};
  \draw[flow] (grid) -- (sim);
  \draw[flow] (sim) -- (pq);
  \draw[flow] (pq) -- (split);
  \draw[flow] (split) -- (win);
  \draw[flow] (win) -- (fit);
  \draw[flow] (fit) -- (cal);
  \draw[flow] (cal) -- (ev);
  \draw[flow] (ev) -- (sel);
\end{tikzpicture}}
\caption{From simulation to a deployable detector. The workbench \ui{Data}
and \ui{Model} tabs run exactly these steps.}
\label{fig:mlpipe}
\end{figure}

\section{Why unsupervised, trained on normal only}

The detectors never see a fault during training. This is not a purity
exercise: a real plant has years of normal operation and a handful of
documented failures, so a supervised classifier would be trained on the rarest
thing it will ever meet and could not recognise a mode nobody labelled.
Instead each detector learns what \emph{normal} looks like and scores how far
a window departs from it. Fault data is used only to measure the result.

Formally, given normal feature vectors $\{\bm x_i\}$ drawn from $p_0$, each
detector defines an anomaly score $a(\bm x)$ that is small where $p_0$ is
dense and large where it is not. A window is flagged when $a(\bm x) > \theta$.

\section{Windows}

A detector looks at a window, not a sample, because what distinguishes faults
are shapes over time: a leak shifts a level, a drift tilts it, stiction makes
it oscillate. One sample cannot express any of those.

\begin{figure}[htbp]
\centering
\scalebox{0.8}{%
\begin{tikzpicture}[x=0.1cm,y=1cm]
  \draw[->] (0,0) -- (125,0) node[right]{sample (1\,Hz)};
  \foreach \s/\c in {0/accent,10/okgreen!70!black,20/warnamber}{
    \draw[\c,very thick,fill=\c!15] (\s,{0.25+\s/40}) rectangle (\s+60,{0.65+\s/40});
    \node[font=\footnotesize,anchor=west] at (\s+61,{0.45+\s/40}) {ends at $\s+59$};
  }
  \draw[<->] (0,-0.25) -- node[below,font=\footnotesize]{window $W=60$} (60,-0.25);
  \draw[<->] (60,-0.25) -- node[below,font=\footnotesize]{stride 10} (70,-0.25);
  \draw[badred,dashed,thick] (75,-0.1) -- (75,1.4) node[above,font=\footnotesize]{fault onset};
\end{tikzpicture}}
\caption{Overlapping windows. Each is labelled by its \emph{last} sample, so a
window straddling the onset is not counted as faulty from its first sample.}
\label{fig:windows}
\end{figure}

Defaults: \code{WINDOW = 60} samples (one minute) and \code{STRIDE = 10}. A
window labelled by its first sample would make detection delay look negative.

\section{Features}

\subsection{Channels}
Nine measured or controller tags are used:
\tg{LT101}, \tg{LT102}, \tg{FT101}, \tg{FT102}, \tg{PT101}, \tg{ZT101},
\tg{FCV101_OP}, \tg{LIC101_SP} and \tg{LIC101_ERR}. Controller output is
included deliberately: in a closed loop the controller absorbs the fault, so
the PV alone hides it (Section~\ref{sec:masking}).

Three more channels are computed from first principles, each because the
fault shows up in a difference that no single tag contains:
\begin{align}
  \text{OP\_MINUS\_ZT} &= \text{FCV101\_OP} - \text{ZT101}, \\
  \text{FLOW\_IMBALANCE} &= \text{FT101} - \text{FT102}, \\
  \text{HEAD\_RESIDUAL} &= \text{PT101} - \frac{\Delta p_0}{1000}
     \Big(1-\big(\tfrac{q}{q_\text{max}}\big)^2\Big),
\end{align}
where $q$ is \tg{FT101} in m$^3$/s.
\tg{OP_MINUS_ZT} is zero on a healthy valve, and a sticking one opens a gap.
\tg{FLOW_IMBALANCE} is zero at steady state by conservation of mass, and a
leak makes it positive: a conservation law, not a learned correlation.
\tg{HEAD_RESIDUAL} is the pump pressure minus what the pump curve says it
should develop at the delivered flow, and impeller wear opens it.
Giving a model the residual of a physical relationship is cheaper and more
robust than making it learn the relationship.

\subsection{Statistics}
Each of the $12$ channels is reduced to four numbers over the window,
$\{\textit{last}, \textit{mean}, \textit{std}, \textit{slope}\}$, giving $48$
features. The slope is the least-squares slope in units per sample,
\begin{equation}
  \beta = \frac{\sum_k (k-\bar k)(x_k - \bar x)}{\sum_k (k-\bar k)^2}.
\end{equation}

\begin{center}
\begin{tabular}{@{}lll@{}}
\toprule
Statistic & Meaning & Where faults show \\
\midrule
last & where it is now & bias, stuck \\
mean & where it has been & leak, fouling, wear \\
std & how much it is moving & stiction (limit cycle), stuck (std $\to$ 0) \\
slope & which way it is going & drift, incipient faults \\
\bottomrule
\end{tabular}
\end{center}

Sixty samples of twelve channels would be 720 nearly redundant numbers. Four
statistics per channel is small enough to train on a modest dataset and keeps
every feature attributable to one channel, which is what makes isolation
possible. Features are named \code{CHANNEL::stat}, e.g.\
\code{FLOW_IMBALANCE::mean}.

\subsection{Standardisation}
Every detector first standardises with statistics of the normal training set:
\begin{equation}
  z_j = \frac{x_j - \mu_j}{\sigma_j}.
\end{equation}
In these units ``normal'' is centred at the origin with unit spread in every
direction, and all scores below are computed on $\bm z$.

\section{Detector 1: z-score limits}

The most straightforward method there is: no model at all. The reconstruction
is the normal mean, which is zero in standardised units, so the score is
\begin{equation}
  a_\text{z}(\bm x) = \sum_{j=1}^{48} z_j^2 .
\end{equation}

\begin{theory}[Interpretation]
If the features were independent Gaussians, $a_\text{z}$ would follow a
$\chi^2$ distribution with 48 degrees of freedom and a fixed limit would set a
known false alarm rate. They are not independent (\tg{LT101} and
\tg{LIC101_ERR} are nearly the same signal), which is why the threshold is
taken empirically from calibration data instead. The z-score is the baseline
every other method must beat; if it cannot be beaten, the honest report says
so.
\end{theory}

\section{Detector 2: PCA and the SPE statistic}

Principal component analysis is the basis of classical multivariate
statistical process control (MSPC). Normal operation of a plant governed by a
few physical relationships occupies a low-dimensional subspace of feature
space; a fault that breaks a relationship pushes data \emph{off} that
subspace.

\subsection{Theory}
With the standardised training matrix $Z\in\mathbb R^{n\times m}$, the
covariance $S = Z^\top Z/(n-1)$ has the eigendecomposition $S = P\Lambda
P^\top$ with eigenvalues $\lambda_1\ge\lambda_2\ge\dots$. Keep the first $k$
eigenvectors $P_k$, chosen as the smallest $k$ that explains 95\% of the
variance:
\begin{equation}
  \frac{\sum_{i\le k}\lambda_i}{\sum_i \lambda_i} \ge 0.95 .
\end{equation}
A window is projected onto the subspace and back, $\hat{\bm z} = P_kP_k^\top
\bm z$, and the score is the \emph{squared prediction error} (SPE, the $Q$
statistic):
\begin{equation}
  a_\text{PCA}(\bm x) = \text{SPE} = \lVert \bm z - \hat{\bm z} \rVert^2
  = \lVert (I - P_kP_k^\top)\,\bm z \rVert^2 .
\end{equation}

\begin{figure}[htbp]
\centering
\scalebox{0.8}{%
\begin{tikzpicture}[scale=1.1]
  \draw[->] (-0.3,0) -- (5,0) node[right]{$z_1$};
  \draw[->] (0,-0.3) -- (0,3.5) node[above]{$z_2$};
  \draw[accent,very thick] (0,0.2) -- (4.8,3.2) node[right]{principal subspace};
  \foreach \x/\y in {0.6/0.7,1.1/0.8,1.5/1.2,1.9/1.3,2.4/1.6,2.8/1.9,3.2/2.2,3.6/2.3,4.0/2.7,1.3/0.9,2.2/1.5,3.0/2.0}
    \fill[slate] (\x,\y) circle (1.4pt);
  \fill[badred] (2.0,2.6) circle (2.2pt) node[above left]{fault};
  \draw[badred,thick,dashed] (2.0,2.6) -- (2.48,1.75) node[midway,right,font=\footnotesize]{$\sqrt{\text{SPE}}$};
\end{tikzpicture}}
\caption{Normal data (grey) lies near the principal subspace. A fault that
breaks the correlation is far from it even when each coordinate is within its
normal range.}
\label{fig:pca}
\end{figure}

The second classical statistic, Hotelling's $T^2 = \bm z^\top P_k
\Lambda_k^{-1} P_k^\top \bm z$, measures unusual variation \emph{within} the
subspace. pyplant uses SPE only, because the faults of interest break
correlations rather than scale them.

\section{Detector 3: isolation forest}

An isolation forest does not model density or reconstruct anything. It exploits
the observation that anomalies are \emph{few and different}, so random
partitioning isolates them in fewer splits than normal points.

\subsection{Theory}
Each of the $t = 200$ trees is built on a random subsample of size $\psi$ by
repeatedly choosing a random feature and a random split value between its
minimum and maximum, until every point is alone. The path length $h(\bm x)$
is the number of splits needed to isolate $\bm x$. Averaged over trees and
normalised by the expected path length of an unsuccessful search in a binary
search tree,
\begin{equation}
  c(\psi) = 2H(\psi-1) - \frac{2(\psi-1)}{\psi}, \qquad
  H(i) \approx \ln i + 0.5772,
\end{equation}
the anomaly score is
\begin{equation}
  a_\text{IF}(\bm x) = 2^{-\mathbb E[h(\bm x)] / c(\psi)} \in (0, 1] .
\end{equation}
Scores near 1 are anomalies, scores well below 0.5 are normal. The code uses
the negated \code{score_samples} of scikit-learn, which is exactly this
quantity, so larger is worse and the threshold and health index work
unchanged.

\begin{figure}[htbp]
\centering
\scalebox{0.8}{%
\begin{tikzpicture}[level distance=9mm,
  level 1/.style={sibling distance=34mm},
  level 2/.style={sibling distance=17mm},
  level 3/.style={sibling distance=9mm},
  every node/.style={circle,draw=slate,fill=panelgrey,minimum size=5mm,inner sep=0pt,font=\scriptsize}]
  \node {}
    child {node[fill=badred!40] {a}}
    child {node {}
      child {node {}
        child {node {n}} child {node {n}}}
      child {node {}
        child {node {n}} child {node {}
          child {node {n}} child {node {n}}}}};
\end{tikzpicture}}
\caption{One isolation tree. The anomaly \textsf{a} is isolated after one
split; normal points \textsf{n} need three or four.}
\label{fig:iforest}
\end{figure}

Because it has no reconstruction, the isolation forest's contributions fall
back to the standardised deviations $z_j^2$ (the same as the z-score
detector).

\section{Thresholds}

Every detector's threshold is a quantile of its score on \emph{held-out
normal} data:
\begin{equation}
  \theta = Q_{q}\big(\{a(\bm x_i) : \bm x_i \in \text{calibration runs}\}\big),
  \qquad q = 0.995 .
\end{equation}
Setting it on the training set measures how well the model memorised, not
how often it will cry wolf. With $q=0.995$, about one normal window in 200
is above threshold; the persistence rule of Chapter~\ref{ch:eval} turns that
into near-zero false alarms.

\section{The health index}

Raw scores have no stable scale: retrain the model and they change entirely.
The health index maps any score onto $0$--$100$ with the threshold pinned at
50 and one decade of score per 50 points:
\begin{equation}
  \text{HI} = \operatorname{clip}\!\Big(50\log_{10}\frac{a}{\theta} + 50,\;0,\;100\Big).
\end{equation}

\begin{figure}[htbp]
\centering
\begin{tikzpicture}
\begin{axis}[width=0.56\textwidth,height=3.8cm,xmode=log,xlabel={score / threshold},
  ylabel={health index},xmin=0.01,xmax=100,ymin=0,ymax=100,grid=major,
  grid style={linegrey}]
\addplot[accent,very thick,domain=0.01:100,samples=200]{min(max(50*log10(x)+50,0),100)};
\addplot[badred,dashed] coordinates {(1,0)(1,100)};
\addplot[badred,dashed] coordinates {(0.01,50)(100,50)};
\node[font=\footnotesize,anchor=west] at (axis cs:1.2,20) {threshold: HI = 50};
\end{axis}
\end{tikzpicture}
\caption{Health index against normalised score. A badly broken plant
(10$\times$ threshold, HI 100) stays distinguishable from a marginal one
(HI just above 50).}
\label{fig:health}
\end{figure}

Pinning the threshold to 50 is what lets an alarm limit on the
\tg{HEALTH_INDEX} tag stay where the alarm engineer put it across a retrain.

\begin{note}[Online inference]
The tag dictionary reserves an \code{ML_RESULT} block (\tg{HEALTH_INDEX},
\tg{TOP_CONTRIB}, \tg{ML_HEARTBEAT}) and the \tg{ALM_HEALTH} alarm for an
inference service that writes the health index back into the control
system. This release computes scores, health index and contributions in
Python and in the workbench \ui{Model} tab; no service writes those tags
yet. Chapter~\ref{ch:extend} sketches how to add one.
\end{note}

\section{Isolation: contributions}

Detection says \emph{that} something is wrong; isolation says \emph{where}.
The reconstruction-based detectors produce a per-feature residual
$\bm r = \bm z - \hat{\bm z}$. Every feature belongs to one channel, so summing
squared residuals by channel and normalising gives each channel's share of the
anomaly:
\begin{equation}
  c_k = \frac{\sum_{j \in \text{channel }k} r_j^2}{\sum_j r_j^2},
  \qquad \sum_k c_k = 1 .
\end{equation}
The \emph{top channel} is $\arg\max_k c_k$. This is the same argument
classical contribution plots rest on. Expected top channels for a fully
developed fault:

\begin{center}
\begin{tabular}{@{}ll@{}}
\toprule
Fault & Expected top channel \\
\midrule
leak & \tg{FLOW_IMBALANCE} \\
pump wear & \tg{HEAD_RESIDUAL} \\
stiction & \tg{OP_MINUS_ZT} \\
fouling & \tg{LT102} or \tg{FT102} \\
sensor faults & not scored (see Chapter~\ref{ch:eval}) \\
\bottomrule
\end{tabular}
\end{center}

\section{Using the detectors from Python}

\begin{lstlisting}
from pyplant.ml.dataset import load, split_scenarios, windowed
from pyplant.ml.detector import make_detector
from pyplant.ml.features import WindowSpec

frame = load("work/demo/dataset.parquet")
data  = windowed(frame, WindowSpec(window=60, stride=10))
split = split_scenarios(frame, seed=0)

det = make_detector("pca", quantile=0.995)     # zscore | pca | isolation_forest | autoencoder
det.fit(data.of_scenarios(split.train).X, data.names)
det.calibrate(data.of_scenarios(split.calibrate).X)

test = data.of_scenarios(split.test_faulty)
hi   = det.health(test.X)          # 0-100, 50 = threshold
top  = det.top_channel(test.X)     # e.g. 'FLOW_IMBALANCE'
det.save("my_detector.joblib")
\end{lstlisting}
